\chapter{Regular Lie groups}
\label{chap:regular}

Regularity asks that the differential equation $\gamma^{-1}\dot\gamma = \xi$ for curves
$\xi$ in the Lie algebra have solutions depending smoothly on $\xi$ \cite[Chapter 6]{GN26}. We
read the logarithmic derivative in the chart $\varphi$ at the identity, and we formalize the
case of continuous curves ($k = 0$ in \cite[Definition 6.1.1]{GN26}), the strongest of the
$C^k$-regularity conditions; the spaces $C^k([0, 1], \mathfrak g)$, $k \geq 1$, are not yet
available as locally convex spaces.

\section{Logarithmic derivatives}

\begin{definition}[Logarithmic derivatives]
  \label{def:log-deriv}
  \lean{LocallyConvex.LieGroups.MBLieGroup.HasLogDerivAt}
  \leanok
  \uses{def:lie-group}
  A curve $\gamma\colon \mathbb{R} \to G$ has logarithmic derivative
  $\delta(\gamma)(t) = \gamma(t)^{-1}\dot\gamma(t) = x$ at $t$ if it is continuous at $t$ and
  $s \mapsto \varphi(\gamma(t)^{-1}\gamma(s))$ has derivative $x$ at $s = t$.
\end{definition}

\begin{lemma}[Uniqueness Lemma]
  \label{lem:log-deriv-unique}
  \lean{LocallyConvex.LieGroups.MBLieGroup.HasLogDerivAt.mul_inv,
    LocallyConvex.LieGroups.MBLieGroup.eq_of_hasLogDerivAt_zero,
    LocallyConvex.LieGroups.MBLieGroup.eq_of_hasLogDerivAt}
  \leanok
  \uses{def:log-deriv}
  Two curves with the same logarithmic derivative that agree at one point are equal
  \cite[Lemma 5.5.4]{GN26}.
\end{lemma}

\begin{proof}
  \leanok
  If $\delta(\gamma_1)(t) = \delta(\gamma_2)(t) = x$, then near $t$,
  $(\gamma_2\gamma_1^{-1})(t)^{-1}(\gamma_2\gamma_1^{-1})(s) = g\, w_2(s) w_1(s)^{-1} g^{-1}$ with
  $g = \gamma_1(t)$ and $w_i(s) = \gamma_i(t)^{-1}\gamma_i(s)$. In the chart this is
  $K(\varphi w_2(s), \varphi w_1(s))$ for a smooth $K$ with $K(p, p) = \varphi(1)$, so its
  derivative is $dK((\varphi(1), \varphi(1)), (x, x)) = 0$: $\gamma_2\gamma_1^{-1}$ has vanishing
  logarithmic derivative. A curve $P$ with vanishing logarithmic derivative is locally constant:
  near $s_0$, $e(s) = \varphi(P(s_0)^{-1}P(s))$ satisfies
  $e(s) = e(s_1) * \varphi(P(s_1)^{-1}P(s))$, so $e' = 0$; and $\mathbb{R}$ is connected.
\end{proof}

\section{$C^0$-regularity}

\begin{definition}[$C^0$-regular Lie groups]
  \label{def:c0-regular}
  \lean{LocallyConvex.LieGroups.MBLieGroup.extendCurve,
    LocallyConvex.LieGroups.MBLieGroup.IsEvol,
    LocallyConvex.LieGroups.MBLieGroup.IsEvol.unique,
    LocallyConvex.LieGroups.MBLieGroup.IsC0Semiregular,
    LocallyConvex.LieGroups.MBLieGroup.evol,
    LocallyConvex.LieGroups.MBLieGroup.IsC0Regular}
  \leanok
  \uses{lem:log-deriv-unique}
  $G$ is $C^0$-semiregular if for every $\xi \in C([0, 1], \mathfrak g)$ there is a curve
  $\gamma_\xi$ with $\gamma_\xi(0) = 1$ and $\delta(\gamma_\xi) = \tilde\xi$, the extension of $\xi$
  to $\mathbb{R}$ by constants; it is unique. $G$ is $C^0$-regular if moreover
  $\mathrm{evol}\colon C([0, 1], \mathfrak g) \to G$, $\xi \mapsto \gamma_\xi(1)$, is smooth
  \cite[Definition 6.1.1]{GN26}.
\end{definition}

\begin{proposition}[Solutions on {$[0, 1]$}]
  \label{prop:regular-interval}
  \lean{LocallyConvex.LieGroups.MBLieGroup.HasLogDerivWithinAt,
    LocallyConvex.LieGroups.MBLieGroup.HasLogDerivWithinAt.hasLogDerivAt_of_Iic_Ici,
    LocallyConvex.LieGroups.MBLieGroup.IsEvolOn,
    LocallyConvex.LieGroups.MBLieGroup.IsEvolOn.oneParam,
    LocallyConvex.LieGroups.MBLieGroup.IsOneParam.isEvol,
    LocallyConvex.LieGroups.MBLieGroup.evol_eq_of_isEvolOn,
    LocallyConvex.LieGroups.MBLieGroup.IsEvolOn.exists_isEvol,
    LocallyConvex.LieGroups.MBLieGroup.IsC0SemiregularGN,
    LocallyConvex.LieGroups.MBLieGroup.isC0Semiregular_iff,
    LocallyConvex.LieGroups.MBLieGroup.evolGN_eq,
    LocallyConvex.LieGroups.MBLieGroup.IsC0RegularGN,
    LocallyConvex.LieGroups.MBLieGroup.isC0Regular_iff}
  \leanok
  \uses{def:c0-regular}
  $C^0$-semiregularity and $C^0$-regularity agree with the notions of
  \cite[Definition 6.1.1]{GN26}, where (6.2) is solved on $[0, 1]$ with one-sided derivatives at
  the endpoints.
\end{proposition}

\begin{proof}
  \leanok
  Restricting gives one direction. Conversely, a solution $c$ on $[0, 1]$ for a constant curve
  $x$ gives the one-parameter group $P(t) = c(1)^{\lfloor t\rfloor} c(t - \lfloor t\rfloor)$: at
  integers, the one-sided logarithmic derivatives are both $x$. A solution $\gamma$ on $[0, 1]$
  for $\xi$ extends by $P_{\xi(0)}$ on $(-\infty, 0]$ and $\gamma(1)P_{\xi(1)}(t - 1)$ on
  $[1, \infty)$, glued at $0$ and $1$ by one-sided logarithmic derivatives. By uniqueness on
  $\mathbb{R}$, every solution on $[0, 1]$ ends at $\mathrm{evol}(\xi)$.
\end{proof}

\begin{theorem}[Regular Lie groups have exponential functions]
  \label{thm:regular-exp}
  \lean{LocallyConvex.LieGroups.MBLieGroup.IsEvol.comp_mul,
    LocallyConvex.LieGroups.MBLieGroup.IsEvol.map_add,
    LocallyConvex.LieGroups.MBLieGroup.IsC0Regular.isExp}
  \leanok
  \uses{def:c0-regular, def:exp}
  If $G$ is $C^0$-regular, $\exp(x) = \mathrm{evol}(x)$ for constant curves $x$ is an exponential
  function of $G$ \cite[Remark 6.1.2(a)]{GN26}.
\end{theorem}

\begin{proof}
  \leanok
  By uniqueness, $\gamma_{cx}(t) = \gamma_x(ct)$, and $\gamma_x(r + a) = \gamma_x(a)\gamma_x(r)$
  since both sides have logarithmic derivative $x$; hence $t \mapsto \exp(tx) = \gamma_x(t)$ is a
  one-parameter group with derivative $x$ at $0$. Smoothness follows from that of $\mathrm{evol}$.
\end{proof}

\begin{proposition}[Vector groups]
  \label{prop:vector-regular}
  \lean{LocallyConvex.LieGroups.continuousAt_of_hasDirDerivAt_one,
    LocallyConvex.LieGroups.VectorGroup.isEvol_integral,
    LocallyConvex.LieGroups.VectorGroup.integralCLM,
    LocallyConvex.LieGroups.VectorGroup.isC0Regular}
  \leanok
  \uses{def:c0-regular, prop:vector-group}
  If $E$ is sequentially complete, $(E, +)$ is $C^0$-regular, with
  $\mathrm{evol}(\xi) = \int_0^1 \xi$ \cite[Proposition 6.1.3]{GN26}.
\end{proposition}

\begin{proof}
  \leanok
  $t \mapsto \int_0^t \tilde\xi$ is the evolution of $\xi$ by the fundamental theorem of calculus.
  The integral is linear, and continuous on $C([0, 1], E)$ because weak integrals of uniformly
  small curves are small.
\end{proof}

\begin{lemma}[Local regularity criterion]
  \label{lem:local-regularity}
  \lean{LocallyConvex.LieGroups.MBLieGroup.HasLogDerivAt.piecewise,
    LocallyConvex.LieGroups.MBLieGroup.IsEvol.concat,
    LocallyConvex.LieGroups.MBLieGroup.evol_eq_mul,
    LocallyConvex.LieGroups.MBLieGroup.exists_mem_valIn,
    LocallyConvex.LieGroups.MBLieGroup.isC0Semiregular_of_local,
    LocallyConvex.LieGroups.MBLieGroup.isC0Regular_of_local,
    LocallyConvex.LieGroups.MBLieGroup.isC0Regular_of_nhds}
  \leanok
  \uses{def:c0-regular}
  If every $\xi$ in an open $0$-neighborhood $U \subseteq C([0, 1], \mathfrak g)$ has an evolution
  and $\mathrm{evol}$ is smooth on $U$, then $G$ is $C^0$-regular \cite[Lemma 6.1.4]{GN26}.
\end{lemma}

\begin{proof}
  \leanok
  $U$ contains $\{\xi \mid \xi([0, 1]) \subseteq W\}$ for an open $0$-neighborhood $W \subseteq E$.
  The halves $\xi_0(t) = \xi(t/2)/2$ and $\xi_1(t) = \xi((1 + t)/2)/2$ of $\xi$ have evolutions
  $\gamma_0$, $\gamma_1$ if $\xi$ has; then $\gamma_0(2t)$ for $t \leq 1/2$ followed by
  $\gamma_0(1)\gamma_1(2t - 1)$ is the evolution of $\xi$, the two pieces being glued at $t = 1/2$,
  where they have the same logarithmic derivative $\xi(1/2)$. Hence
  $\mathrm{evol}(\xi) = \mathrm{evol}(\xi_0)\,\mathrm{evol}(\xi_1)$. If $\xi$ takes values in
  $2^m W$, its halves take values in $2^{m-1} W$; by induction on $m$, evolutions exist and
  $\mathrm{evol}$ is smooth on the open sets $\{\xi \mid \xi([0, 1]) \subseteq 2^m W\}$, which cover
  $C([0, 1], \mathfrak g)$ since the image of each $\xi$ is compact, hence bounded. Glöckner and
  Neeb cut $\xi$ into $N$ pieces at once.
\end{proof}

